% ============================================================
% INFLUENCE OF DESIGN PARAMETERS
% ============================================================
\section*{Influence of Design Parameters on Aircraft Wing Performance}

The evolution of airfoil design, 
as documented in seminal works such as NACA TR-610 \cite{NACA_TR_610}, 
demonstrates a systematic approach to correlating geometric definitions with aerodynamic performance. 
By examining the fundamental NACA families, 
ranging from the early 4-digit series to the laminar-flow 6-series, 
designers can tailor aircraft performance through specific parametric adjustments.

% ----------------------------------------------------------
% Classical Thickness Distributions
% ----------------------------------------------------------

\subsection*{The NACA 4-Digit Series: Geometric Fundamentals}

The NACA 4-digit series utilizes three geometric parameters: maximum camber, camber position, and maximum thickness. 
Increasing the maximum camber enhances the curvature of the mean camber line, 
which directly raises the lift coefficient at a zero angle of attack and increases the magnitude of the negative pitching moment about the quarter-chord. 
While thin airfoil theory suggests that camber primarily shifts the lift curve upward without altering its slope, 
excessive camber can induce adverse pressure gradients, 
leading to premature boundary layer separation.

% ----- Representative Profiles -----

\subsubsection*{Representative Profiles}
NACA 0012 (Symmetrical standard for stabilizers), 
NACA 2412 (Commonly used in general aviation wings like the Cessna 172), 
and NACA 4415 (High-lift variant for heavy utility aircraft).

% ----------------------------------------------------------
% Classical Thickness Distributions
% ----------------------------------------------------------

\subsection*{The NACA 5-Digit Series: Optimized Loading}

The NACA 5-digit series represents a shift toward theoretical aerodynamic loading. 
Unlike the 4-digit series, the camber line is derived to meet a specific design lift coefficient, with the maximum camber typically positioned further forward. 
This forward shift significantly enhances maximum lift capabilities while maintaining lower pitching moments compared to 4-digit counterparts. 
A unique feature of this series is the reflex indication (the third digit), where a reflexed camber line, 
characterized by upward curvature at the trailing edge, 
is employed to neutralize nose-down pitching tendencies, a critical requirement for tailless or flying-wing configurations.

% ----- Representative Profiles -----

\subsubsection*{Representative Profiles}
NACA 23012 (Renowned for its high lift-to-moment ratio, used on the B-17 and Beechcraft King Air) 
and NACA 23112 (Reflexed version for specialized stability).

% ----------------------------------------------------------
% Classical Thickness Distributions
% ----------------------------------------------------------

\subsection*{Modified 4- and 5-Digit Series: High-Speed Refinement}

The modified series introduces additional control over the leading-edge radius and the chordwise position of maximum thickness. 
Adjusting the leading-edge radius allows designers to balance stall characteristics against high-speed efficiency; a larger radius reduces sensitivity to angle-of-attack variations and delays separation, 
whereas a sharper radius may optimize specific drag regimes. 
Moving the position of maximum thickness, 
traditionally at 30 percent of the chord, 
further aft (to 40 percent or 50 percent) is a primary strategy 
for delaying drag divergence and improving high-subsonic performance, 
though it often necessitates a trade-off in structural spar depth.

% ----- Representative Profiles -----

\subsubsection*{Representative Profiles}
NACA 0009-64 (Used in high-speed jet tail surfaces) 
and NACA 23012-64 (Modified for improved Mach number performance in early jet designs like the F-84).

% ----------------------------------------------------------
% Classical Thickness Distributions
% ----------------------------------------------------------

\subsection*{NACA 1-Series and 16-Series: High-Speed Specialization}

While the 1-series was the precursor to laminar flow research, 
the 16-series became the benchmark for high-speed, low-pressure applications. These airfoils are designed to maintain a nearly uniform pressure distribution over the chord, 
effectively eliminating the high-suction peaks that lead to flow separation or cavitation. 
Their geometry is characterized by a very small leading-edge radius and maximum thickness located near 50 percent of the chord, 
making them ideal for high-speed propellers and hydrofoils.

% ----- Representative Profiles -----

\subsubsection*{Representative Profiles}
NACA 16-012 (Standard for high-speed propeller sections) 
and NACA 16-212 (Cambered version for racing aircraft).

% ----------------------------------------------------------
% Classical Thickness Distributions
% ----------------------------------------------------------

\subsection*{NACA 6 and 6A Series: Laminar Flow Innovation}

The NACA 6-series airfoils are engineered to maximize the extent of the laminar boundary layer by maintaining a favorable pressure gradient over a large portion of the chord. 
This results in the "drag bucket" phenomenon, a specific range of lift coefficients where profile drag is exceptionally low. 
The design lift coefficient dictates the center of this bucket, 
while thickness influences both the width of the bucket and the structural integrity of the wing. 
The 6A series represents a practical modification where the trailing-edge "cusp" is removed 
to simplify manufacturing and improve pitching moment behavior without significant loss of laminar efficiency.

% ----- Representative Profiles -----

\subsubsection*{Representative Profiles}
NACA 64-212 (Classic laminar flow wing) 
and NACA 64A210 (Modernized version for high-performance jets like the F-16).
